% ================================================
% CARTOON EXAMPLE - Hebrew Academic Template
% ================================================
% This example demonstrates \cartoon{} - the RTL side-illustration wrap column
% introduced in v7.4.0. Compiles to approximately 4-5 pages.
% Copyright (c) 2025 Dr. Segal Yoram. All rights reserved.
%
% NOTE: this file deliberately loads NO extra packages and defines NO commands.
% Everything below is a class built-in. \cartoon needs no \usepackage at all.

\documentclass{hebrew-academic-template}

% Image directory for \cartoon (note the trailing slash)
\setcartoonpath{images/cartoons/}

% Default wrap column width, applied book-wide unless overridden per call
\setcartoonwidth{0.32\textwidth}

\hebrewtitle{דוגמת איורי צד - תבנית אקדמית עברית}
\englishtitle{Cartoon Wrap Column Example - Hebrew Academic Template \clsversion}
\hebrewauthor{ד"ר סגל יורם}
\date{\textenglish{July 2026}}

\begin{document}

\maketitle
\tableofcontents
\listoffigures
\newpage

% ==================== INTRODUCTION ====================

\hebrewsection{מהו איור צד: \entoc{What Is a Side Illustration}}

הפקודה \en{\textbackslash cartoon} ממקמת איור קטן בעמודת עטיפה בצד תחילת הקריאה
של גוש הטקסט, כלומר בצד ימין במסמך עברי. הטקסט העברי זורם לצד האיור, וברגע
שהוא עובר את גובה האיור הוא חוזר לרוחב המלא של גוש הטקסט. התוצאה דומה לספר
מאויר שבו האיורים יושבים בשוליים ולא קוטעים את רצף הקריאה.

זהו ההבדל המרכזי בין \en{\textbackslash cartoon} לבין \en{\textbackslash hebrewfigure}:

\begin{itemize}
    \item \en{\textbackslash hebrewfigure} יוצר איור צף לרוחב מלא, שקוטע את הטקסט.
    \item \en{\textbackslash cartoon} יוצר עמודת עטיפה צרה, שהטקסט זורם סביבה.
\end{itemize}

השימוש המיועד הוא איורים קטנים, קריקטורות והמחשות המלוות את הטקסט. לדיאגרמות,
גרפים וטבלאות יש להמשיך להשתמש ב-\en{\textbackslash hebrewfigure}.

% ==================== BASIC USAGE ====================

\hebrewsection{שימוש בסיסי: \entoc{Basic Usage}}

התחביר הבסיסי מקבל שלושה ארגומנטים: שם הקובץ, הכיתוב בעברית, והתווית להפניה.

\begin{english}
\begin{verbatim}
\cartoon{cart_person}{כיתוב האיור בעברית}{fig:label}
\end{verbatim}
\end{english}

הפסקה הבאה מדגימה את התוצאה. שימו לב שקריאת האיור יושבת בשורה נפרדת בין
פסקאות, מיד לפני הפסקה שאמורה לזרום סביבו.

\cartoon{cart_person}{דמות פשוטה: איור הצד יושב בצד ימין של גוש הטקסט,
והטקסט זורם לצידו}{fig:cart-person}
כאן מתחיל הטקסט שזורם לצד האיור. האיור ממוקם בצד ימין, שהוא צד תחילת הקריאה
בעברית, והטקסט נכרך סביבו באופן טבעי. ברגע שהטקסט עובר את גובה האיור הוא
חוזר לרוחב המלא של גוש הטקסט וממשיך עד השוליים הימניים. כפי שניתן לראות,
הכיתוב מופיע מתחת לאיור, ממוספר אוטומטית, ונכלל ברשימת האיורים שבתחילת
המסמך. ההפניה לאיור נעשית כרגיל: איור~\ref{fig:cart-person}. שורה נוספת של
טקסט להדגמת הזרימה. שורה נוספת של טקסט. שורה נוספת של טקסט. שורה נוספת של
טקסט. שורה נוספת של טקסט. שורה נוספת של טקסט. שורה נוספת של טקסט.

הפסקה הזו מגיעה אחרי האיור והיא לרוחב מלא לחלוטין.

% ==================== PLACEMENT RULES ====================

\hebrewsection{כללי מיקום: \entoc{Placement Rules}}

זהו החלק החשוב ביותר במסמך זה. עמודת עטיפה רגישה למיקום, והמחלקה אינה יכולה
לאכוף את הכללים האלה עבורכם.

\hebrewsubsection{הכלל העיקרי: בין פסקאות: \entoc{Between Paragraphs}}

קריאת \en{\textbackslash cartoon} חייבת לשבת \textbf{בין פסקאות}, בשורה
נפרדת, מיד לפני הפסקה שאמורה לזרום סביב האיור:

\begin{english}
\begin{verbatim}
...סוף הפסקה הקודמת.

\cartoon{file}{caption}{fig:label}
הפסקה שנכרכת סביב האיור מתחילה כאן...
\end{verbatim}
\end{english}

זו דרישה מתועדת של מנגנון העטיפה. קריאה \textbf{בתוך} פסקה עשויה להיראות
תקינה במקרה, אך אינה נתמכת ועלולה להישבר באופן בלתי צפוי.

\hebrewsubsection{מה להימנע ממנו: \entoc{What to Avoid}}

\begin{enumerate}
    \item \textbf{מיד אחרי כותרת:} אל תמקמו איור צד מיד אחרי
          \en{\textbackslash hebrewsection} או \en{\textbackslash hebrewsubsection}.
    \item \textbf{בתוך פסקה:} ראו הכלל העיקרי לעיל.
    \item \textbf{ליד רשימה:} רשימות ועמודת עטיפה אינן משתלבות היטב.
    \item \textbf{ליד איור צף אחר:} שני מנגנוני מיקום המתחרים על אותו שטח.
    \item \textbf{שני איורי צד באותה פסקה:} השתמשו בפסקה נפרדת לכל אחד.
\end{enumerate}

% ==================== CONFIGURATION ====================

\hebrewsection{הגדרות: \entoc{Configuration}}

\hebrewsubsection{נתיב התמונות: \entoc{Image Path}}

הפקודה \en{\textbackslash setcartoonpath} קובעת את תיקיית התמונות, כך שקריאות
האיור מקבלות שם קובץ בלבד. שימו לב ללוכסן בסוף:

\begin{english}
\begin{verbatim}
\setcartoonpath{images/cartoons/}
\end{verbatim}
\end{english}

\hebrewsubsection{רוחב ברירת המחדל: \entoc{Default Width}}

הפקודה \en{\textbackslash setcartoonwidth} קובעת את רוחב עמודת העטיפה עבור כל
הספר. שינוי במקום אחד מכוונן מחדש את כל האיורים:

\begin{english}
\begin{verbatim}
\setcartoonwidth{0.32\textwidth}
\end{verbatim}
\end{english}

\hebrewsubsection{דריסת רוחב נקודתית: \entoc{Per-Call Width Override}}

ארגומנט אופציונלי ראשון דורס את רוחב ברירת המחדל עבור קריאה בודדת:

\begin{english}
\begin{verbatim}
\cartoon[0.42\textwidth]{file}{caption}{fig:label}
\end{verbatim}
\end{english}

\cartoon[0.42\textwidth]{cart_clock}{שעון: איור זה משתמש ברוחב מותאם של
\num{0.42} מרוחב הטקסט במקום ברירת המחדל}{fig:cart-clock}
הפסקה הזו זורמת סביב איור רחב יותר. עמודת העטיפה תופסת כאן חלק גדול יותר
מרוחב גוש הטקסט, ולכן נותרות פחות מילים בכל שורה. זהו איזון שיש לשקול: איור
רחב יותר קריא יותר, אך הטקסט לצידו נעשה צר ומקוטע. שורה נוספת של טקסט. שורה
נוספת של טקסט. שורה נוספת של טקסט. שורה נוספת של טקסט. שורה נוספת של טקסט.
שורה נוספת של טקסט. שורה נוספת של טקסט. שורה נוספת של טקסט.

ככלל, רוחב שבין \num{0.30} ל-\num{0.40} מרוחב הטקסט עובד היטב לרוב האיורים.

% ==================== PAGE BREAK GUARD ====================

\hebrewsection{שובר העמוד האוטומטי: \entoc{Automatic Page-Break Guard}}

איור צד הנפתח קרוב לתחתית העמוד עלול לגלוש מעבר לשוליים התחתונים ולהידפס על
גבי הכותרת התחתונה. זו תקלה מוכרת של מנגנוני עטיפה.

הפקודה \en{\textbackslash cartoon} מטפלת בכך אוטומטית: היא מודדת את התמונה
בגודלה הסופי, מוסיפה קצבה לכיתוב, ובודקת אם נותר די מקום בעמוד. אם לא, העמוד
נשבר והאיור עובר לעמוד הבא בשלמותו.

אין צורך לעשות דבר כדי להפעיל את ההגנה הזו. האיור הבא ממוקם בכוונה קרוב
לתחתית העמוד, וההגנה מעבירה אותו לעמוד הבא במקום לאפשר לו לגלוש.

\cartoon{cart_books}{ערימת ספרים: איור זה ממוקם בכוונה קרוב לתחתית העמוד כדי
להדגים את שובר העמוד האוטומטי}{fig:cart-books}
אם ההגנה פועלת, האיור הזה נמצא כעת בראש עמוד חדש ואינו חוצה את הכותרת
התחתונה. שורה נוספת של טקסט. שורה נוספת של טקסט. שורה נוספת של טקסט. שורה
נוספת של טקסט. שורה נוספת של טקסט. שורה נוספת של טקסט. שורה נוספת של טקסט.
שורה נוספת של טקסט. שורה נוספת של טקסט. שורה נוספת של טקסט.

שימו לב שההגנה מטפלת בגלישה מתחתית העמוד בלבד. היא אינה יכולה למנוע התנגשות
מסוג אחר: כותרת סעיף שנוחתת בתוך עמודת עטיפה שעדיין פעילה. אם הטקסט שאחרי
האיור קצר מדי, הכותרת הבאה תיכנס אל תוך שטח העטיפה והכיתוב יתנגש בה.

לכן הכלל המעשי הוא לוודא שאחרי כל איור צד יש די טקסט כדי לסגור את עמודת
העטיפה לפני הכותרת הבאה. פסקה זו קיימת בדיוק לשם כך, והיא מדגימה את הכמות
המינימלית הסבירה. שורה נוספת של טקסט. שורה נוספת של טקסט. שורה נוספת של
טקסט. שורה נוספת של טקסט. שורה נוספת של טקסט. שורה נוספת של טקסט.

% ==================== PREPARING IMAGES ====================

\hebrewsection{הכנת קובצי האיור: \entoc{Preparing Image Files}}

\hebrewsubsection{יחס גובה-רוחב: \entoc{Aspect Ratio}}

זהו הכלל החשוב ביותר בהכנת התמונות: \textbf{התמונה חייבת להיות גבוהה מרוחבה},
ביחס שבין \num{1}:\num{1} ל-\num{4}:\num{5}.

הסיבה היא חשבונית. עמודת עטיפה ברוחב \num{0.32} מגוש טקסט של כ-\num{14}
סנטימטרים היא כ-\num{4.1} סנטימטרים. תמונה בפורמט רחב של \num{16}:\num{9}
תהיה אז בגובה \num{2.3} סנטימטרים בלבד. תמונה לאורך ביחס \num{4}:\num{5}
תהיה בגובה \num{5.1} סנטימטרים - יותר מכפול שטח ציור.

האיור הבא מדגים את הבעיה. זהו איור רחב, ובגודל ההדפסה בפועל הפרטים בו כמעט
בלתי קריאים:

\cartoon{cart_supplied}{דוגמה נגדית: איור בפורמט רחב. שימו לב עד כמה הוא נמוך
בעמודת העטיפה ועד כמה קשה להבחין בפרטיו}{fig:cart-landscape}
כפי שניתן לראות, איור רחב מתכווץ לרצועה נמוכה. סצנה בת שלוש דמויות הופכת
לכתם בלתי קריא, ועמודת העטיפה נעשית רדודה מכדי שהטקסט יזרום סביבה באופן
משמעותי. חשוב להדגיש: \textbf{חיתוך אינו פותר את הבעיה}. גזירת החלק העליון של
תמונה רחבה רק מגדילה את יחס הרוחב לגובה ומחמירה את המצב.

האיור הזה מדגים גם את הכלל שהוזכר בסעיף הקודם: יש להשאיר די טקסט אחרי האיור
כדי שעמודת העטיפה תיסגר לפני הכותרת הבאה. שורה נוספת של טקסט. שורה נוספת של
טקסט. שורה נוספת של טקסט. שורה נוספת של טקסט. שורה נוספת של טקסט. שורה נוספת
של טקסט. שורה נוספת של טקסט. שורה נוספת של טקסט.

\hebrewsubsection{ללא טקסט בתוך התמונה: \entoc{No Text Inside the Image}}

התמונה עצמה לא תכיל טקסט, כותרת או רצועת כיתוב. הכיתוב מגיע מפקודת
\en{\textbackslash caption} של המחלקה, ולכן טקסט המוטבע בתמונה יופיע פעמיים:
פעם כפיקסלים בגופן זר וגודל קבוע, ופעם כטקסט אמיתי.

טקסט המוטבע בתמונה גם אינו ניתן לעריכה, אינו ניתן לחיפוש, ואינו נכלל ברשימת
האיורים.

\hebrewsubsection{פורמט וגודל: \entoc{Format and Size}}

\begin{itemize}
    \item \textbf{\en{PNG}} לאיורי קו וקריקטורות - דחיסה ללא אובדן.
    \item \textbf{\en{PDF}} לגרפים ודיאגרמות וקטוריות.
    \item רזולוציה של לפחות \num{1200} פיקסלים בצלע הארוכה.
    \item רקע לבן או שקוף.
\end{itemize}

% ==================== BEST PRACTICES ====================

\hebrewsection{עקרונות מומלצים: \entoc{Best Practices}}

\begin{enumerate}
    \item \textbf{מיקום בין פסקאות} - לעולם לא בתוך פסקה ולא אחרי כותרת.
    \item \textbf{תמונה לאורך} - גבוהה מרוחבה, אחרת האיור נעלם בגודל ההדפסה.
    \item \textbf{ללא טקסט בתמונה} - הכיתוב מגיע מ-\en{\textbackslash caption}.
    \item \textbf{כיתוב מסביר} - מה רואים באיור ומה לומדים ממנו.
    \item \textbf{הפניה בגוף הטקסט} - כל איור מוזכר באמצעות \en{\textbackslash ref}.
    \item \textbf{רוחב אחיד} - הגדירו רוחב אחד ב-\en{\textbackslash setcartoonwidth}
          ודרסו נקודתית רק בעת הצורך.
    \item \textbf{לא יותר מדי} - איור צד אחד לכל כמה עמודים. צפיפות פוגעת בקריאות.
\end{enumerate}

% ==================== SUMMARY ====================

\hebrewsection{סיכום: \entoc{Summary}}

מסמך זה הדגים:

\begin{itemize}
    \item שימוש בסיסי ב-\en{\textbackslash cartoon} עם שלושה ארגומנטים
    \item כלל המיקום בין פסקאות ומה להימנע ממנו
    \item הגדרת נתיב באמצעות \en{\textbackslash setcartoonpath}
    \item הגדרת רוחב באמצעות \en{\textbackslash setcartoonwidth} ודריסה נקודתית
    \item שובר העמוד האוטומטי המונע גלישה מעבר לשוליים
    \item דרישות התמונה: פורמט לאורך, ללא טקסט מוטבע
    \item הפניות צולבות ורשימת איורים
\end{itemize}

איורי הצד שהוצגו: איור~\ref{fig:cart-person}, איור~\ref{fig:cart-clock},
איור~\ref{fig:cart-books} ואיור~\ref{fig:cart-landscape}.

\end{document}
